\section{Introduction}
Digital twin applications are built upon computational models that imitate physical systems and rely on robust handling of real-time data. However, the high complexity in coupling these models often leads to monolithic integration solutions. As a result many 'traditional' digital twins are designed around large, singular models that generate outputs strictly from predefined inputs. These systems usually enclose an entire physical process, covering structural behavior, fluid interactions, control logic and more, within a single, all-encompassing simulation. While this approach can capture multiple phenomena, it may reduce the digital twin to static mapping of in- and outputs, offering only limited insight beyond the immediate, single-purpose functionality. Consequently, such designs lack the flexibility to be updated or scaled selectively when individual modeling components change, and they do not readily support deeper interaction among various analytical methods or predictive algorithms.

These limitations become increasingly apparent as physical systems become more complex. First, scalability poses a challenge: a single simulation environment can become computationally expensive if more detailed analyses or higher fidelity models are introduced. Attempting to expand one part of a monolithic model often requires re-engineering the entire system to maintain consistency, blocking iterative improvements. Second, modularity may be compromised when the entire workflow, from input pre-processing to output generation, lives under a unified architecture. Replacing or upgrading one aspect risks disrupting the whole simulation and demands extensive refactoring. Third, resilience suffers when all processes depend one another such that failure in one sub-component can propagate rapidly, taking down the entire digital twin. Furthermore, because a single model typically prioritizes overarching system objectives, interaction among sub-models may not receive the nuanced treatment each domain requires.

These gaps motivate the exploration of a distributed multi-agent approach, wherein the digital twin does not serve as just an input-output intermediary but an ecosystem of specialized, cooperating entities. By loosely coupling several autonomous agents, each designed to handle distinct tasks such as structural simulation, fluid behavior analysis, or predictive maintenance, the digital twin can evolve beyond a unidirectional mapping of inputs to outputs. Instead, it supports bidirectional or even network-wide data exchanges, enabling domain-specific optimizations and real-time adaptability.

In this architecture, scalability is enhanced because each agent operates independently; if one simulation component requires more resources, additional instances of that agent can be assigned in distributed computing environments without reworking the entire system. Modularity emerges from the fact that agents communicate through well-defined interfaces or middleware, upgrading or swapping out a specific model affects only that agent and its interfaces, leaving the rest of the digital twin intact. This multi-agent system also brings robustness and resilience, as a failure or delay in one agent need not trigger cascading failures across the whole platform; other agents can continue functioning as long as data flows remain valid and consistent.

Adopting a multi-agent digital twin introduces its own set of research challenges. Designing an agent-based approach requires careful consideration of how agents coordinate, share data, and synchronize individual computations. Moreover, identifying the best practices from existing Multi-Agent Systems design principles, such as communication protocols, interaction patterns, and failure strategies, is critical to ensure robust, reliable operation in (near)real-time environments. Technology choices also become essential: selecting tools, platforms, and protocols that can effectively facilitate agent orchestration and data management is essential for practical deployments. Finally, demonstrating the feasibility and tangible benefits of this approach calls for the development of a prototype system that showcases how a digital twin, represented as a multi-agent system, may outperform or supplements monolithic designs.


Despite these challenges, a distributed, multi-agent digital twin may deliver more than an input-output pipeline. Each agent not only processes specialized data, but also learns, adapts, and collaborates, providing a richer, more flexible representation of the physical system. This paves the way for the integration of algorithms, utilizing new computational paradigms, and offering deeper insight into the interactions between system components. Thus, the multi-agent approach aspires to transform the digital twin from a static mirror of reality into a dynamic, continually evolving partner in system analysis, optimization, and decision-making.


\section*{Research Question}
\textbf{Main Question:}\\
How can digital twin applications be modularized and integrated using a multi-agent architecture?

\noindent \textbf{Sub-questions}
\begin{enumerate}
    \item \textbf{RQ1} What are the limitations of current approaches to Digital Twin modularization and integration? \\
    \item \textbf{RQ2} What opportunities do Multi-Agent Systems offer to address these gaps? \\
    \item \textbf{RQ3} What are the current best practices in Multi-Agent System architecture design principles that can be applied to DT modularization and integration? \\
    \item \textbf{RQ4} What technologies, tools, and protocols are available for implementing agent-based Digital Twin systems? \\
    \item \textbf{RQ5} How can we develop a prototype to demonstrate the feasibility and benefits of a multi-agent architecture for Digital Twins? \\
\end{enumerate}
